\chapter{Koszul duality reference tables}
\label{app:koszul-reference}


\begin{remark}[Relationship to \S\ref{app:koszul_higher_genus}]
The genus-graded Koszul duality theory is developed in
\S\ref{app:koszul_higher_genus}. This appendix provides reference
tables and the essential image characterization; the core definitions
and theorems are stated there and should not be duplicated. Where a
result appears in both locations, the version in
\S\ref{app:koszul_higher_genus} is canonical.
\end{remark}

\section{Genus-graded Koszul duality}

\begin{theorem}[Extended Koszul duality; \ClaimStatusConditional]
\label{thm:extended-koszul-appendix}
Let $(\mathcal{A},\mathcal{A}^{!})$ be a genus-$0$ chiral Koszul
pair and assume the higher-genus bar objects are finite-window
complete, mapping-class equivariant, and centrally curved in the
sense of Theorem~\ref{thm:higher-genus-inversion}. Put
\[
\mathcal{A}^{\mathrm i,(g)}:=\barB^{\mathrm{ch},(g)}(\mathcal{A}),
\qquad
\mathcal{A}^{!,(g)}:=
\mathbb{D}_{\Ran}\mathcal{A}^{\mathrm i,(g)}
\]
on the Verdier-finite lane. Then the completed genus series
\[
\widehat{\mathcal{A}}=\prod_{g\geq 0}\hbar^{2g-2}\mathcal{A}^{(g)},
\qquad
\widehat{\mathcal{A}^{\mathrm i}}
=\prod_{g\geq 0}\hbar^{2g-2}\mathcal{A}^{\mathrm i,(g)}
\]
is a compatible higher-genus Koszul system: the counits
\[
\Omega^{\mathrm{ch}}_g
\mathcal{A}^{\mathrm i,(g)}
\xrightarrow{\;\sim\;}
\mathcal{A}^{(g)}
\]
are quasi-isomorphisms on the higher-genus Koszul locus, and
Verdier duality sends the bar-dual coalgebra
$\mathcal{A}^{\mathrm i,(g)}$ to the algebra
$\mathcal{A}^{!,(g)}$.
\end{theorem}

\begin{proof}
This is Theorem~\ref{thm:extended-koszul} in
\S\ref{app:koszul_higher_genus}, proved via the Prism Principle
(Theorem~\ref{thm:prism-higher-genus}), strict nilpotence at all genera
(Theorem~\ref{thm:genus-induction-strict}), and the universal extension tower
(Theorem~\ref{thm:universal-extension-tower}).
\end{proof}

\section{Definition and basic properties}

\begin{definition}[Genus-graded Koszul algebra]
A genus-graded associative algebra $\mathcal{A} = \bigoplus_{g \geq 0} \mathcal{A}^{(g)}$ is \emph{Koszul} if:
\[\text{Ext}_{\mathcal{A}^{(g)}}^{i,j}(\mathbb{k}, \mathbb{k}) = 0 \text{ for } i \neq j\]
where the bigrading is by homological degree and internal degree, and the Koszul property holds at each genus.
\end{definition}

\begin{remark}[Classical vs chiral Koszulness]\label{rem:classical-vs-chiral-koszulness}
The criterion above is \emph{classical} algebraic Koszulness
(BGS~\cite{BGS96}). In the chiral setting, classical Koszulness of
the enveloping algebra does \emph{not} imply chiral Koszulness of the
corresponding vertex algebra, because the chiral bar construction
uses all OPE poles (Borcherds locality), not just the leading term.

\emph{Concrete illustration.}
The enveloping algebra $U(\mathfrak{sl}_2)$ is Koszul in the classical
sense: $\operatorname{Ext}^{i,j}_{U(\mathfrak{sl}_2)}(\mathbb{k},
\mathbb{k}) = 0$ for $i \neq j$ (this follows from the PBW theorem
and the Koszul complex for the symmetric algebra). For the affine
algebra $\widehat{\mathfrak{sl}}_2$ at level~$k$, the OPE of
currents has both a simple pole (the structure constants $f^{ab}_c$)
and a \emph{double} pole (the level term $k\langle a, b\rangle$).
The classical bar construction sees only the simple pole and factors
through the Chevalley--Eilenberg complex of~$\mathfrak{sl}_2$. The
chiral bar construction sees the full OPE, including the double pole,
so additional cancellations must be verified: specifically, the PBW
spectral sequence (Theorem~\ref{thm:pbw-koszulness-criterion}) must
degenerate, which requires controlling the interaction of double-pole
contributions across all bar degrees. When it does degenerate,
$\widehat{\mathfrak{sl}}_{2,k}$ is chiral Koszul, but this is a
\emph{theorem} (proved in Chapter~\ref{chap:kac-moody-koszul}), not a
corollary of the classical statement.

See Definition~\ref{def:chiral-koszul-pair} for the chiral Koszul pair
formulation.
\end{remark}

\begin{remark}[Quadratic candidate dual versus full chiral Koszulness:
MA-12 endpoint discipline]
\label{rem:quadratic-vs-full-koszulness-koszul-ref}
\index{MA-12 (finite-spin / quadratic / classical = full theorem)}
The Gui--Li--Zeng quadratic chiral duality~\cite{GLZ22} produces a
\emph{candidate} Koszul dual $\cA^!_{\mathrm{GLZ}} = \cA(s^{-1}N^\vee
\omega^{-1}, P^\perp)$ together with a Maurer--Cartan comparison map
\[
\operatorname{Hom}_{\mathrm{ChirAlg}}(\cA, \cB)
\;\longleftrightarrow\;
\MC\bigl(\cA^!_{\mathrm{GLZ}} \otimes \cB\bigr).
\]
This is \emph{not} full chiral Koszulness. Bare promotion of the
quadratic candidate to the full Koszulness theorem is the master
pattern MA-12 (\emph{quadratic = full theorem}); the endpoint
hypothesis package required to lift the candidate to the
inversion-theorem statement
$\Omega^{\mathrm{ch}}\Bbar^{\mathrm{ch}}(\cA) \xrightarrow{\sim} \cA$
has three independent ingredients (a) quadratic-dual candidate;
(b) MC comparison map upgraded to a quasi-isomorphism of convolution
dg Lie algebras; (c) bar--cobar adjunction $K^2 \simeq \id$ on the
comparison locus. The full statement is
Theorem~\ref{thm:full-koszulness-separate} in
Appendix~\ref{app:dual-methodology}; this remark is its
forward-reference inside the genus-graded reference apparatus.

Type signatures separating the two:
\emph{Quadratic candidate} carries (Open quadrant; quadratic chiral
data; level $1\!\rightarrow\!2$; HP = $E_\infty$-chiral + quadratic
$(N,P)$ + dualizable). \emph{Full Koszulness} carries (Open quadrant;
bar--cobar inversion; level $1\!\leftrightarrow\!2$; HP = quadratic
candidate + MC quasi-isomorphism + $K^2 \simeq \id$ on comparison
locus). The first is necessary, the second is sufficient.
\end{remark}

\begin{theorem}[Genus-graded Koszul duality theorem; \ClaimStatusConditional]
\label{thm:genus-graded-koszul-duality-appendix}
Let $\mathcal{A}$ be genus-graded Koszul, with each
$\mathcal{A}^{(g)}$ connected, augmented, finite-window complete, and
centrally curved. The chiral bar-dual coalgebra is
\[
\mathcal{A}^{\mathrm i}:=
\bigoplus_{g\geq 0}\barB^{\mathrm{ch},(g)}(\mathcal{A}),
\]
completed $\hbar$-adically when infinitely many genera contribute.
On the Verdier-finite lane the Koszul dual algebra is
\[
\mathcal{A}^{!}:=\mathbb{D}_{\Ran}(\mathcal{A}^{\mathrm i}).
\]
The classical formula
\[
\bigoplus_{g\geq 0}
\operatorname{Ext}^{*}_{\mathcal{A}^{(g)}}(\mathbb{k},\mathbb{k})
\]
computes the same algebra only on the ordinary quadratic finite-type
locus, where classical linear duality identifies Ext with the dual of
the bar-dual coalgebra. Outside that locus it computes a classical
associated-graded shadow — a quadratic Koszul candidate, in the
sense of Gui--Li--Zeng~\cite{GLZ22}, that requires the additional
MC-comparison and bar--cobar inversion ingredients of
Remark~\ref{rem:quadratic-vs-full-koszulness-koszul-ref} and
Theorem~\ref{thm:full-koszulness-separate} to lift to a full
Koszulness statement. For genuine vertex algebras the object
$\mathcal{A}^{\mathrm i}$ is the chiral bar coalgebra, and
$\mathcal{A}^{!}$ is the post-Verdier algebra
$\mathbb{D}_{\Ran}\mathcal{A}^{\mathrm i}$.
\end{theorem}

\begin{proof}
This is Theorem~\ref{thm:genus-graded-koszul-duality} in
\S\ref{app:koszul_higher_genus}, proved by applying the $\Eone$-chiral Koszul
duality (Theorem~\ref{thm:e1-chiral-koszul-duality}) genus by genus, with
nilpotence ensured by Theorem~\ref{thm:genus-induction-strict} and
involutivity by bar-cobar inversion
(Theorem~\ref{thm:bar-cobar-inversion-qi}).
\end{proof}

\subsection{Genus-graded chiral Koszul duality}

For chiral algebras across all genera the data must remember the
coalgebra before Verdier duality.

\begin{definition}[Genus-graded chiral Koszul duality]
A genus-graded chiral Koszul datum is a system
\[
(\mathcal{A}^{(g)},\mathcal{C}^{(g)},\tau^{(g)})_{g\geq 0}
\]
where $\mathcal{A}^{(g)}$ is an augmented chiral algebra,
$\mathcal{C}^{(g)}$ is a coaugmented conilpotent chiral coalgebra
or a finite-window topologically conilpotent completion,
and $\tau^{(g)}\colon\mathcal{C}^{(g)}\to\mathcal{A}^{(g)}$ is a
chiral twisting morphism. It is Koszul if
\[
\mathcal{A}^{(g)}\otimes_{\tau^{(g)}}\mathcal{C}^{(g)}
\simeq \mathbb{C}
\]
in the derived category of chiral modules for every genus~$g$, and if
the finite-window system is equivariant for the mapping-class group.
Then
\[
\mathcal{A}^{\mathrm i,(g)}=\mathcal{C}^{(g)},
\qquad
\mathcal{A}^{!,(g)}=\mathbb{D}_{\Ran}\mathcal{C}^{(g)}
\]
whenever the Verdier dual algebra is defined.
\end{definition}

\subsection{Curved and filtered generalizations across genera}

\begin{definition}[Genus-graded curved Koszul duality]
A genus-graded curved $A_\infty$ chiral algebra
$(\mathcal{A}^{(g)},m_0^{(g)},m_1^{(g)},m_2^{(g)},\ldots)$ satisfies
\[
(m_1^{(g)})^2=[m_0^{(g)},-].
\]
Its bar-dual object is the curved chiral coalgebra
\[
\mathcal{A}^{\mathrm i,(g)}=\barB^{\mathrm{ch},(g)}(\mathcal{A}),
\]
not the algebra $\mathcal{A}^{!,(g)}$. When the finite-window Verdier
duality hypotheses hold, the dual algebra is
$\mathcal{A}^{!,(g)}=\mathbb{D}_{\Ran}\mathcal{A}^{\mathrm i,(g)}$.
The dual curvature is transported by the suspended Verdier pairing;
the sign is the suspension sign fixed in
Appendix~\ref{app:signs}. Mapping-class covariance is a
condition on the whole family over the universal curve, not on a
single period matrix.
\end{definition}

\subsection{Computational tools across genera}

\begin{lemma}[Genus-graded Koszul complex resolution; \ClaimStatusConditional]
\label{lem:genus-graded-koszul-resolution-appendix}
For genus-graded Koszul $\mathcal{A}$ with central curvature, the
minimal chiral resolution of $\mathbb{k}$ at genus~$g$ is the twisted
bar resolution
\[
\cdots \to
\mathcal{A}^{(g)}\otimes
(\mathcal{A}^{\mathrm i,(g)})_{(2)}
\to
\mathcal{A}^{(g)}\otimes
(\mathcal{A}^{\mathrm i,(g)})_{(1)}
\to
\mathcal{A}^{(g)}
\to
\mathbb{k}.
\]
The differential is induced by the universal twisting morphism
$\tau^{(g)}\colon\mathcal{A}^{\mathrm i,(g)}\to\mathcal{A}^{(g)}$.
In the geometric model its coefficients are logarithmic propagators
on $\overline{C}_n(\Sigma_g)$, together with the OPE residue maps
along collision strata. If finite-type Verdier duality is available,
dualizing the coalgebra factors rewrites the same resolution in terms
of the algebra $\mathcal{A}^{!,(g)}$.
\end{lemma}

\begin{proof}
This is Lemma~\ref{lem:genus-graded-koszul-resolution} in
\S\ref{app:koszul_higher_genus}, proved by combining the
geometric bar construction
(Theorem~\ref{thm:geometric-equals-operadic-bar}) with the higher-genus
bar-cobar inversion (Theorem~\ref{thm:higher-genus-inversion}).
\end{proof}

\subsection{Physical interpretation across genera}

The genus-graded Koszul system supplies the algebraic skeleton for the
genus expansion of string amplitudes
(Chapter~\ref{ch:genus-expansions}). Electric-magnetic duality,
open-closed duality, and mirror symmetry enter only through the
connection theorems that name their hypotheses.

\subsection{Genus-graded Maurer--Cartan elements and twisting}

\begin{theorem}[Genus-graded MC elements parametrize deformations; \ClaimStatusProvedElsewhere]
\label{thm:genus-graded-mc-appendix}
Let
\[
\mathfrak{g}^{(g)}_{\mathcal{A}}
=\operatorname{Conv}\bigl(\mathcal{A}^{\mathrm i,(g)},\mathcal{A}^{(g)}\bigr)
\]
be the completed convolution dg Lie algebra controlling genus-$g$
chiral operations. Formal genus-$g$ deformations are Maurer--Cartan
elements
\[
\alpha^{(g)}\in
\mathfrak{g}^{(g),1}_{\mathcal{A}}[[t]],
\qquad
D^{(g)}\alpha^{(g)}
\;+\;\frac{1}{2}[\alpha^{(g)},\alpha^{(g)}]=0.
\]
For a fixed curved background $\Theta^{(g)}$ this is equivalently the
perturbation equation
\[
D_{\Theta^{(g)}}\beta^{(g)}
\;+\;\frac{1}{2}[\beta^{(g)},\beta^{(g)}]=0.
\]
The twisted operator
$D^{(g)}_{\alpha}=D^{(g)}+[\alpha^{(g)},-]$ squares to zero exactly
when the displayed MC equation holds. Gauge equivalence is the
Deligne--Getzler--Hinich groupoid of the pronilpotent convolution
algebra. Logarithmic one-form representatives and classical
Yang--Baxter equations occur on the current/r-matrix sublocus; they
are not the definition of a general chiral MC element.
\end{theorem}

\begin{proof}
This is Theorem~\ref{thm:genus-graded-mc} in
\S\ref{app:koszul_higher_genus}. The controlling dg Lie algebra is
the chiral Hochschild convolution algebra
(Theorem~\ref{thm:linfty-cochains}), realized geometrically by
kernels on configuration spaces
(Theorem~\ref{thm:geometric-chiral-hochschild}). The equality
$(D+[\alpha,-])^2=0$ is the standard dg Lie computation and is
equivalent to the MC equation. The higher-genus compatibility follows
from mapping-class equivariance of the bar differential
(Theorem~\ref{thm:prism-higher-genus}).
\end{proof}

\subsection{Koszul duality at higher genus: the tower structure}
\label{app:koszul-reference-higher-genus}

The genus-$0$ counit
\[
\Omega^{\mathrm{ch}}\mathcal{A}^{\mathrm i,(0)}
\xrightarrow{\;\sim\;}
\mathcal{A}^{(0)}
\]
extends to higher genus only on the higher-genus Koszul locus and in
the completed curved ambient of Theorem~\ref{thm:higher-genus-inversion}.

\subsubsection{\texorpdfstring{The genus $g$ statement}{The genus g statement}}

For each $g\geq 0$ in that locus there is a counit
\[
\Omega^{(g)}\mathcal{A}^{\mathrm i,(g)}
\xrightarrow{\;\sim\;}
\mathcal{A}^{(g)}.
\]
Here $\Omega^{(g)}$ is the genus-$g$ cobar construction and
$\mathcal{A}^{\mathrm i,(g)}$ is the genus-$g$ bar-dual coalgebra.

\subsubsection{Compatibility}

The boundary maps are residue maps along the Deligne--Mumford
boundary. A nonseparating boundary component lowers the genus by one;
a separating boundary component lands in a tensor product of the two
lower-genus factors. Thus the tower is not a single chain
$C^{(g)}\to C^{(g-1)}$. It is a modular-operadic system whose arrows
are the clutching and residue maps
\[
\operatorname{Res}_{\mathrm{nonsep}}\colon
\mathcal{A}^{\mathrm i,(g)}
\to \mathcal{A}^{\mathrm i,(g-1)},\qquad
\operatorname{Res}_{\mathrm{sep}}\colon
\mathcal{A}^{\mathrm i,(g)}
\to
\mathcal{A}^{\mathrm i,(g_1)}\widehat{\otimes}
\mathcal{A}^{\mathrm i,(g_2)}.
\]
The cobar counits commute with these maps by the modular operad
axioms.

\subsubsection{The limit}

The all-genus object is the $\hbar$-adic completed product
\[
\widehat{\mathcal{A}}
=\prod_{g\geq 0}\hbar^{2g-2}\mathcal{A}^{(g)}.
\]
Finite-window quotients, not the raw direct sum over all genera, carry
the strict algebraic topology used by Theorem~\ref{thm:higher-genus-inversion}.

\subsubsection{Modular invariance}

At each genus $g$, the counit is equivariant for the mapping-class
group $\Gamma_g=\operatorname{MCG}(\Sigma_g)$:
\[
\Omega^{(g)}(\sigma^*\mathcal{A}^{\mathrm i,(g)})
\simeq
\sigma^*\mathcal{A}^{(g)}
\qquad(\sigma\in\Gamma_g).
\]
Scalar genus-$g$ amplitudes obtained by integration over the proper
Deligne--Mumford stack are therefore modular invariant when the input
chain is mapping-class equivariant.


\section{Classification of chiral algebras by Koszul type}
\label{app:koszul-classification}

See Proposition~\ref{prop:classification-table} in Appendix~\ref{app:existence-criteria} for the complete classification table of chiral algebras by Koszul type; see \S\ref{sec:filtered-vs-curved-comprehensive} for the underlying theory.

%================================================================
% SECTION: ESSENTIAL IMAGE CHARACTERIZATION
%================================================================

\section{\texorpdfstring{Essential image: when is $\widehat{\mathcal{C}} = \mathcal{A}^{\mathrm i}$?}{Essential image: when is C = A i?}}
\label{sec:essential-image-koszul}

\subsection{The characterization problem}

\begin{question}[Inverse problem]\label{q:inverse-problem-koszul}
Given a chiral coalgebra $\widehat{\mathcal{C}}$, when does there
exist a chiral algebra $\mathcal{A}$ such that
\[
\widehat{\mathcal{C}}\simeq\mathcal{A}^{\mathrm i}
=\barB^{\mathrm{ch}}(\mathcal{A})?
\]
When this holds, the Koszul dual algebra is
$\mathcal{A}^{!}=\mathbb{D}_{\Ran}\widehat{\mathcal{C}}$ on the
Verdier-finite lane.

The invariant form of the question asks for the \emph{essential image}
of the chiral bar coalgebra functor.
\end{question}

\begin{remark}\label{rem:why-essential-image-matters}
The essential image question has four parts: recognition of bar-dual
coalgebras, completeness of the bar functor, uniqueness of
$\mathcal{A}$ up to the bar-cobar equivalence, and reconstruction by
$\Omega^{\mathrm{ch}}\widehat{\mathcal{C}}$.
\end{remark}

\subsection{Main characterization theorem}

\begin{theorem}[Essential image of the chiral bar-dual coalgebra; \ClaimStatusProvedHere]\label{thm:essential-image-koszul}
A coaugmented chiral coalgebra $\widehat{\mathcal{C}}$ in the
finite-type complete $E_1$-chiral lane is equivalent to the bar-dual
coalgebra $\mathcal{A}^{\mathrm i}$ of a chiral algebra
$\mathcal{A}$ if and only if the following five conditions hold.

\begin{enumerate}
\item \emph{Coaugmented connected.} The coradical is the ground field,
and the coaugmentation and counit split it:
 \[\mathbb{C}\xrightarrow{\eta}\widehat{\mathcal{C}}
 \xrightarrow{\epsilon}\mathbb{C},
 \qquad \epsilon\eta=\operatorname{id}_{\mathbb{C}}.\]

\item \emph{Algebraic or topological conilpotency.} In the algebraic
lane every element of $\ker(\epsilon)$ is killed by a high enough
iterate of the reduced coproduct:
\[
\ker(\epsilon)=\bigcup_{n\geq 0}\ker(\bar{\Delta}^{(n)}).
\]
In the completed lane there are finite-window quotients
$\widehat{\mathcal{C}}_{\leq N}$ such that
$\widehat{\mathcal{C}}\simeq\varprojlim_N
\widehat{\mathcal{C}}_{\leq N}$, each
$\widehat{\mathcal{C}}_{\leq N}$ is conilpotent, and the reduced
coproduct is continuous and lowers the window filtration.

\item \emph{Geometric representability.} $\widehat{\mathcal{C}}$
admits a presentation as a factorization coalgebra on $\Ran(X)$,
computed by logarithmic forms on the Fulton--MacPherson
compactifications and compatible with collision residues.

\item \emph{Curvature compatibility.} In the strict lane the curvature
acts centrally, so the internal differential squares to zero. In the
curved lane the identity is the CDG relation $d^2=[\mu_0,-]$ and the
object is read in the coderived/contraderived Positselski ambient.

\item \emph{Bar-cobar unit.} The cobar algebra
$\Omega^{\mathrm{ch}}(\widehat{\mathcal{C}})$ exists in the same
ambient, and the adjunction unit
\[
\eta_{\widehat{\mathcal{C}}}\colon
\widehat{\mathcal{C}}\longrightarrow
\barB^{\mathrm{ch}}\Omega^{\mathrm{ch}}(\widehat{\mathcal{C}})
\]
is a quasi-isomorphism in the strict finite-type lane, respectively a
coacyclic equivalence in the curved finite-window lane.
\end{enumerate}

When these conditions hold,
\[
\mathcal{A}=\Omega^{\mathrm{ch}}(\widehat{\mathcal{C}})
\]
is the reconstructed algebra and
$\widehat{\mathcal{C}}\simeq\mathcal{A}^{\mathrm i}$. On the
Verdier-finite lane the Koszul dual algebra is
$\mathcal{A}^{!}=\mathbb{D}_{\Ran}\widehat{\mathcal{C}}$. The object
$\mathcal{A}^{!}$ is not the cobar algebra
$\Omega^{\mathrm{ch}}(\widehat{\mathcal{C}})$; the latter is the
bar-cobar reconstruction of $\mathcal{A}$.
\end{theorem}

\begin{proof}
Necessity is the construction of the chiral bar coalgebra. The
coaugmentation is dual to the augmentation of $\mathcal{A}$, the
reduced coproduct is deconcatenation and lowers tensor length, the
Fulton--MacPherson model is Definition~\ref{def:geometric-bar}, and
curvature satisfies the curved $A_\infty$ or CDG identity. The
bar-cobar unit is the coalgebra-side form of
Theorem~\ref{thm:e1-chiral-koszul-duality} in the strict lane and of
Theorem~\ref{thm:chiral-positselski-weight-completed} in the completed
curved lane.

Conversely, condition (5) gives
\[
\widehat{\mathcal{C}}\simeq
\barB^{\mathrm{ch}}\Omega^{\mathrm{ch}}(\widehat{\mathcal{C}})
\]
in the stated ambient. With
$\mathcal{A}:=\Omega^{\mathrm{ch}}(\widehat{\mathcal{C}})$ this says
exactly that $\widehat{\mathcal{C}}$ is the bar-dual coalgebra
$\mathcal{A}^{\mathrm i}$. Applying $\mathbb{D}_{\Ran}$, when the
Verdier-finiteness hypotheses hold, produces the algebra
$\mathcal{A}^{!}$ and not a second cobar reconstruction.
\end{proof}

\subsection{Conilpotency and connectedness}

\begin{lemma}[Algebraic and completed conilpotency; \ClaimStatusProvedHere]\label{lem:conilpotency-necessary}
If
$\widehat{\mathcal{C}}\simeq
\mathcal{A}^{\mathrm i}=\barB^{\mathrm{ch}}(\mathcal{A})$
for some augmented chiral algebra $\mathcal{A}$, then the algebraic
bar coalgebra is conilpotent. Its finite-window completion is
topologically conilpotent, not necessarily elementwise conilpotent.
\end{lemma}

\begin{proof}
Let $\bar{\mathcal{A}}$ be the augmentation ideal. The underlying
coalgebra of the chiral bar construction is the ordered tensor
coalgebra
\[
T^c(s^{-1}\bar{\mathcal{A}})
=\bigoplus_{n\geq 0}(s^{-1}\bar{\mathcal{A}})^{\otimes n}
\]
with the deconcatenation coproduct, realized geometrically on
ordered configuration spaces. If $c$ has tensor length at most $N$,
then the $(N+1)$-fold reduced coproduct $\bar\Delta^{(N+1)}(c)$
vanishes. Every element in the algebraic bar coalgebra has finite
tensor length. Hence
\[
\ker(\epsilon)=\bigcup_{N\geq 0}\ker(\bar\Delta^{(N)}).
\]
In the completed finite-window model an element may have infinitely
many nonzero tensor-length components. The correct assertion is
therefore
\[
\widehat{\mathcal{C}}
\simeq
\varprojlim_N \widehat{\mathcal{C}}_{\leq N},
\]
where each quotient $\widehat{\mathcal{C}}_{\leq N}$ is conilpotent
and the reduced coproduct is continuous for the inverse-limit
topology. No fixed iterate of $\bar\Delta$ need kill an arbitrary
completed element before passing to a finite window.
\end{proof}

\begin{lemma}[Connectedness characterizes augmentation; \ClaimStatusProvedHere]\label{lem:connectedness-augmentation}
For finite-type objects, coaugmented connected coalgebras are dual to
augmented algebras.
\end{lemma}

\begin{proof}
If $C$ is coaugmented and connected, the splitting
$\mathbb{C}\xrightarrow{\eta} C\xrightarrow{\epsilon}\mathbb{C}$
dualizes to an algebra splitting
\[
\mathbb{C}\xrightarrow{u} C^\vee\xrightarrow{\varepsilon}\mathbb{C}.
\]
Thus $C^\vee$ is augmented. Conversely, the augmentation of a
finite-type algebra $A$ dualizes to a coaugmentation of $A^\vee$, and
the unit of $A$ dualizes to the counit of $A^\vee$. Connectedness is
the assertion that the coradical is exactly the image of the
coaugmentation.
\end{proof}

\subsection{Geometric representability}

\begin{definition}[Geometrically representable coalgebra]\label{def:geom-representable-coalgebra}
A chiral coalgebra $\widehat{\mathcal{C}}$ is \emph{geometrically representable}
if there exists:
\begin{enumerate}
\item A factorization algebra $\mathcal{F}$ on configuration spaces $\{C_n(X)\}$
\item A quasi-isomorphism:
 \[\widehat{\mathcal{C}} \simeq \int_{C_\bullet(X)} \mathcal{F}\]
 (factorization homology)
\end{enumerate}
\end{definition}

\begin{theorem}[Bar-dual coalgebras are geometrically representable; \ClaimStatusProvedHere]\label{thm:koszul-geom-rep}
If
$\widehat{\mathcal{C}}\simeq
\mathcal{A}^{\mathrm i}=\barB^{\mathrm{ch}}(\mathcal{A})$
for a chiral algebra $\mathcal{A}$, then $\widehat{\mathcal{C}}$ is
geometrically representable via
\[\mathcal{A}^{\mathrm i} \simeq \bar{B}^{\text{geom}}(\mathcal{A}) =
\bigoplus_{n \geq 0} \Gamma\left(\overline{C}_n(X), \mathcal{A}^{\boxtimes n}
\otimes \Omega^\bullet\right)\]
\end{theorem}

\begin{proof}
The geometric bar complex (Definition~\ref{def:geometric-bar}) provides the
geometric realization:

The factorization algebra is
\[\mathcal{F}_{C_n}(U) = \Gamma(U, \mathcal{A}^{\boxtimes n}|_U)\]
for $U \subseteq C_n(X)$.

Factorization homology is computed by the bar complex
\[
\bar{B}^{\text{geom}}(\mathcal{A})
= \int_{C_\bullet(X)} \mathcal{F}
\]
by Theorem~\ref{thm:bar-factorization-homology}.

By Theorem~\ref{thm:bar-concentration} on the Koszul locus, the
geometric bar coalgebra is the bar-dual coalgebra
$\mathcal{A}^{\mathrm i}$. The algebra $\mathcal{A}^{!}$ appears only
after Verdier duality:
\[
\mathcal{A}^{!}=\mathbb{D}_{\Ran}\mathcal{A}^{\mathrm i}.
\]
\end{proof}

\begin{corollary}[Converse with the unit criterion; \ClaimStatusConditional]\label{cor:geom-implies-koszul}
If $\widehat{\mathcal{C}}$ is geometrically representable by a factorization algebra
$\mathcal{F}$ on configuration spaces, is coaugmented connected,
algebraically or topologically conilpotent, curvature-compatible, and
finite-window complete, and if the unit
\[
\eta_{\widehat{\mathcal{C}}}\colon
\widehat{\mathcal{C}}\to
\barB^{\mathrm{ch}}\Omega^{\mathrm{ch}}(\widehat{\mathcal{C}})
\]
is a quasi-isomorphism in the strict lane or a coacyclic equivalence
in the curved completed lane, then
\[\widehat{\mathcal{C}} \simeq \mathcal{A}^{\mathrm i}\]
for $\mathcal{A} = \Omega^{\mathrm{ch}}(\widehat{\mathcal{C}})$.
\end{corollary}

\begin{proof}
Define $\mathcal{A}=\Omega^{\mathrm{ch}}(\widehat{\mathcal{C}})$.
The connected and conilpotent hypotheses make the reduced coproduct
available and make the cobar differential convergent, either
elementwise in the algebraic lane or continuously in finite windows.
In the strict lane $d_\Omega^2=0$ by coassociativity and compatibility
with the internal differential. In the curved lane its square is the
prescribed inner curvature, and the object is interpreted in the
coderived Positselski ambient.

By Theorem~\ref{thm:geometric-equals-operadic-bar}, the geometric bar
construction of $\mathcal{A}$ agrees with the operadic bar construction.
The unit hypothesis gives
\[
\widehat{\mathcal{C}}\simeq
\barB^{\mathrm{ch}}\Omega^{\mathrm{ch}}(\widehat{\mathcal{C}})
=\barB^{\mathrm{ch}}(\mathcal{A})
\]
as chiral coalgebras, strictly or in the coderived category according
to the lane. Thus
$\widehat{\mathcal{C}}\simeq\mathcal{A}^{\mathrm i}$. The geometric
representability identifies this coalgebra with the corresponding
factorization-homology model on configuration spaces. On the
Verdier-finite lane $\mathbb{D}_{\Ran}$ gives the algebra
$\mathcal{A}^{!}$.
\end{proof}

\subsection{Curvature and centrality}

\begin{theorem}[Central curvature and internal strictness; \ClaimStatusProvedHere]\label{thm:curvature-central-appendix}
Let $\mathcal{A}$ be a curved $A_\infty$ chiral algebra with
zero-ary operation $m_0$. The unary operator $m_1$ is an internal
differential if and only if $m_0$ acts centrally:
\[
m_1^2=[m_0,-]=0.
\]
The stronger uncurved condition is $m_0=0$. If $m_0$ is not central,
the correct object is curved/CDG and must be read in the coderived or
completed curved lane.
\end{theorem}

\begin{proof}
We verify centrality directly from the curved $A_\infty$ relations
(cf.~Positselski~\cite{Positselski11}).

A curved $A_\infty$-algebra carries two operators with different
square laws:
\begin{itemize}
\item The \emph{bar differential} $d_{\bar{B}}$ on the tensor coalgebra
$\bar{B}(\mathcal{A}) = (T^c(s^{-1}\bar{\mathcal{A}}),\, d_{\bar{B}})$,
which encodes all the operations $m_0, m_1, m_2, \ldots$
simultaneously. In a curvature-killed or period-corrected total model
this operator is square-zero. In the raw curved CDG model its square is
the inner curvature action.
\item The \emph{internal differential} $m_1$ on the algebra
$\mathcal{A}$ itself. In the curved case, $m_1^2 \neq 0$ in general:
$m_1$ fails to be a differential, and the failure is measured by the
curvature~$m_0$.
\end{itemize}

The degree-$0$ relation gives the cocycle condition.
The curved $A_\infty$ relation at degree~$n=0$ has a single term
$(r,s,t) = (0,0,0)$, giving $m_1(m_0) = 0$. Thus $m_0$ is a cocycle
for~$m_1$.

The degree-$1$ relation gives the curvature identity.
At degree~$n=1$, the three terms $(r,s,t)$ with $r+s+t=1$ are:
\begin{align*}
(0,1,0) &: \; (-1)^{0 \cdot 1 + 0}\, m_1(m_1(a)) = m_1^2(a), \\
(0,0,1) &: \; (-1)^{0 \cdot 0 + 1}\, m_2(m_0, a) = -m_2(m_0, a), \\
(1,0,0) &: \; (-1)^{1 \cdot 0 + 0}\, m_2(a, m_0) = m_2(a, m_0).
\end{align*}
Setting the sum to zero:
\begin{equation}\label{eq:curvature-identity}
m_1^2(a) = m_2(m_0, a) - m_2(a, m_0) = [m_0, a]_{m_2}
\quad \text{for all } a.
\end{equation}
This is the \emph{curvature identity}: $m_1$ squares not to zero but
to the inner derivation $[m_0, -]$.

From~\eqref{eq:curvature-identity}, two regimes emerge:

\begin{enumerate}[label=(\roman*)]
\item \emph{Curved} ($m_0$ not central): $[m_0, a]_{m_2} \neq 0$ for
some~$a$, so $m_1^2 \neq 0$, and $\mathcal{A}$ has no internal
differential. The bar differential $d_{\bar{B}}$ on $\bar{B}(\mathcal{A})$
is then a curved/CDG differential with square governed by the same
inner curvature. The completed bar object is well-defined in the
curved Positselski ambient, not as a raw strict dg coalgebra.

\item \emph{Internally strict} ($m_0$ central): $[m_0, a]_{m_2} = 0$
for all~$a$, so $m_1^2 = 0$ and $m_1$ is an honest differential.
The algebra $(\mathcal{A}, m_1)$ is then a genuine dg algebra. It is
uncurved only in the stronger case $m_0=0$.
\end{enumerate}

The standard examples fix the normalization.
\begin{itemize}
\item \emph{Heisenberg at level~$\kappa$.} The curvature is
$m_0 = -\kappa \cdot \omega$, a scalar multiple of the symplectic
pairing. Scalars commute with everything:
$[-\kappa \omega, a]_{m_2} = 0$ for all~$a$. So $\mathcal{H}_\kappa$
has an internally strict differential at every level. It is uncurved
only at $\kappa=0$.

\item \emph{Kac--Moody at level~$k$.} The shifted curvature is
proportional to $k+h^\vee$. The completed bar object away from the
critical level is curved and satisfies the CDG identity
$d_{\barB,k}^{\,2}=[m_0,-]$. Strict square-zero occurs at the critical
level or after imposing a twisted flat model in which the curvature
has been killed. The scalar invariant remains
\[
\kappa(V_k(\mathfrak{g}))
=\frac{\dim(\mathfrak{g})(k+h^\vee)}{2h^\vee}.
\]
\end{itemize}

In the bar construction $\bar{B}(\mathcal{A})$, the curvature $\mu_0$
is the image of $m_0$ under the desuspension
$s^{-1}\colon\mathcal{A}\to s^{-1}\mathcal{A}$. The curved
$A_\infty$ identities encode the two displayed relations:
$m_1(m_0)=0$ and $m_1^2=[m_0,-]$. Centrality of $m_0$ is precisely the
extra condition under which the internal differential itself squares
to zero; vanishing of $m_0$ is the separate condition that removes
curvature altogether.
\end{proof}

\begin{example}[Virasoro central charge]\label{ex:virasoro-central-charge-curvature}
The Virasoro central extension is determined by the Lie algebra
$2$-cocycle $\frac{1}{12}(m^3-m)\delta_{m+n,0}$. The scalar central
charge commutes with all $L_n$, but the completed Virasoro bar model
still carries a model-dependent curved term whose scalar invariant is
\[
\kappa(\Vir_c)=c/2.
\]
Strictness is a statement about the chosen curved model, not a
consequence of the centrality of the symbol $c$ alone.
\end{example}

\subsection{Formal completeness}

\begin{definition}[Coaugmentation coideal]\label{def:coaugmentation-coideal}
For a coaugmented connected coalgebra
$\mathbb{C}\xrightarrow{\eta}\widehat{\mathcal{C}}
\xrightarrow{\epsilon}\mathbb{C}$, the \emph{coaugmentation coideal}
is
\[
\bar{\mathcal{C}}=\ker(\epsilon)\simeq
\operatorname{coker}(\eta).
\]
It is the reduced coalgebra on which the reduced coproduct
$\bar\Delta$ is defined. In the completed lane the finite-window
topology is the inverse-limit topology coming from the quotients that
retain bounded tensor length, not the naive powers of a multiplication
ideal.
\end{definition}

\begin{theorem}[Completion characterization; \ClaimStatusConditional]\label{thm:completion-characterization}
\textup{(Positselski \cite{Positselski11})}
A coalgebra $\widehat{\mathcal{C}}$ satisfying the geometric,
connected, conilpotent, and curvature-compatible hypotheses of
Theorem~\ref{thm:essential-image-koszul} can be a completed
bar-dual coalgebra $\mathcal{A}^{\mathrm i}$ only in the
finite-window topology:
\[
\widehat{\mathcal{C}}
\simeq
\varprojlim_N \widehat{\mathcal{C}}_{\leq N},
\]
where $\widehat{\mathcal{C}}_{\leq N}$ is the quotient keeping bar
length at most $N$. Conversely, such a complete coalgebra is
bar-dual to
$\mathcal{A}:=\Omega^{\mathrm{ch}}(\widehat{\mathcal{C}})$ exactly
when the unit criterion of Theorem~\ref{thm:essential-image-koszul}
holds.
\end{theorem}

\begin{proof}
\emph{($\Rightarrow$) Necessity.}
If $\widehat{\mathcal{C}}\simeq\barB^{\mathrm{ch}}(\mathcal{A})$,
the tensor-length filtration gives finite bar windows
$\widehat{\mathcal{C}}_{\leq N}$. The completed bar coalgebra is the
inverse limit of these windows by construction.

\emph{($\Leftarrow$) Sufficiency.}
If $\widehat{\mathcal{C}}$ is complete in this topology and satisfies
the unit criterion, define
\[
\mathcal{A}=\Omega^{\mathrm{ch}}(\widehat{\mathcal{C}}).
\]
Completeness gives convergence of the cobar construction, while the
unit criterion gives
\[
\barB^{\mathrm{ch}}(\mathcal{A}) \simeq \widehat{\mathcal{C}}
\]
in the strict or coderived ambient specified there.

Therefore
$\widehat{\mathcal{C}}\simeq\mathcal{A}^{\mathrm i}$, and
$\mathcal{A}^{!}$ is obtained from
$\widehat{\mathcal{C}}$ by Verdier duality when that duality is
available.
\end{proof}

\subsection{Uniqueness of the algebra}

\begin{theorem}[Uniqueness up to quasi-isomorphism; \ClaimStatusConditional]\label{thm:uniqueness-algebra}
If $\widehat{\mathcal{C}}\simeq \mathcal{A}^{\mathrm i}
\simeq \mathcal{B}^{\mathrm i}$ for two chiral algebras
$\mathcal{A}$ and $\mathcal{B}$, then:
\[\mathcal{A} \simeq \mathcal{B}\]
(quasi-isomorphic as chiral algebras).
\end{theorem}

\begin{proof}
The cobar construction provides canonical algebra reconstructions:
\[
\mathcal{A}\simeq
\Omega^{\mathrm{ch}}(\mathcal{A}^{\mathrm i})
\simeq
\Omega^{\mathrm{ch}}(\widehat{\mathcal{C}})
\simeq
\Omega^{\mathrm{ch}}(\mathcal{B}^{\mathrm i})
\simeq
\mathcal{B}.
\]

These quasi-isomorphisms are bar-cobar inversion maps
(Theorem~\ref{thm:bar-cobar-inversion-qi}); they do not identify the
coalgebra $\widehat{\mathcal{C}}$ with either algebraic Koszul dual
$\mathcal{A}^{!}$ or $\mathcal{B}^{!}$ before Verdier duality.
\end{proof}

\begin{remark}[Non-uniqueness at the strict level]\label{rem:non-uniqueness-strict}
The theorem only guarantees quasi-isomorphism, not strict isomorphism. Different
presentations of the same chiral algebra (e.g., different choices of generators
and relations) give strictly different algebras that are quasi-isomorphic.

\emph{Example.} The Heisenberg chiral algebra can be presented after
choosing different Darboux bases of the same symplectic vector space:
\begin{itemize}
\item generators $(a_i,a_i^*)$ with OPE
$a_i(z)a_j^*(w)\sim k\delta_{ij}/(z-w)^2$;
\item generators $(x_i,p_i)$ obtained from $(a_i,a_i^*)$ by a
symplectic change of basis.
\end{itemize}

These are different presentations, but they define isomorphic chiral
algebras and hence the same bar-dual coalgebra.
\end{remark}
